\section*{Data}

I would use a quarterly panel of US-domiciled, US-listed manufacturing firms (SIC 2000 to 3999) with December fiscal year ends and verifiable US production locations. Manufacturing provides a direct link to physical assets and their replacement. The sample would cover eight pre-storm quarters (2022Q3 to 2024Q2) and four post-storm quarters (2024Q4 to 2025Q3). Two prior annual cycles would help assess investment trends, while the following year would capture initial responses. I would exclude 2024Q3 so that each observation covers a quarter entirely before or after Helene.

Financial data would come from Compustat North America through WRDS, subject to obtaining access. Company annual reports available through the SEC would provide production locations reported before Helene. The Federal Emergency Management Agency (FEMA) supplies disaster declarations identifying counties approved for assistance after the storm.

Compustat's \texttt{CAPXY} records capital expenditure accumulated within the fiscal year. I would calculate quarterly capex by subtracting the previous quarter's accumulated amount within that year. The first quarter needs no subtraction. Investment ($I_{it}$) would equal 100 times quarterly capex divided by total assets at the end of 2023. This fixed denominator prevents the ratio from rising simply because reported assets fall after the storm. Cash holdings ($C_i$) would equal cash and short-term investments (\texttt{CHEQ}) divided by total assets (\texttt{ATQ}), both measured at the end of 2023, before Helene could affect the balance sheet.

The exposure dummy ($E_i$) would equal one if at least one principal US production facility is in a county designated for FEMA Individual Assistance following Helene, and zero otherwise. This measures geographic exposure rather than confirmed firm damage. Firms without such facilities would form the control group. I would check whether both groups include comparable industries and have similar pre-storm investment patterns. Production locations are used because headquarters may be elsewhere.

I would require positive assets at the end of 2023 and exclude firms with incomplete location records or overlapping facility exposure to other major storms in 2024. I would check mergers and divestitures for changes in reporting scope. Firms leaving the market would remain for available periods, and missing observations would not be replaced by zero.
